\chapter{Testing and Evaluation}
\label{ch:testing}

\section{Introduction}
This chapter reports how the system was verified and how well it performs. It
covers three kinds of evidence. \emph{Software testing} establishes that the code
does what it was written to do, with particular attention to the properties that
would otherwise fail silently. \emph{Model evaluation} establishes how accurate
the forecasts are, reported per horizon and under rolling validation rather than
as a single number. \emph{System validation} establishes that the deployed
service actually works, by exercising every published endpoint against the live
deployment. The chapter closes with the operational evidence from the running
system, an assessment of the threats to the validity of these results, and a
discussion of what the results mean.

\section{Testing Objectives}
Five objectives were set.

\begin{enumerate}
  \item To verify that the index computation is correct at its breakpoints, its
  category boundaries and its edge cases, since every number the system produces
  is expressed in that unit.
  \item To verify that no target leakage exists in the feature representation,
  because leakage produces a model that scores well and fails in production, and
  the failure is invisible in the metrics that would normally be trusted.
  \item To establish that the model has genuine forecasting skill by comparing it
  against a persistence baseline at every horizon, not merely on average.
  \item To estimate forward performance honestly, using rolling validation rather
  than a single fortunate split.
  \item To confirm that the deployed system serves every documented capability,
  and that capabilities which are unavailable fail in the specified, informative
  way rather than obscurely.
\end{enumerate}

\section{Testing Strategy}
Testing was applied at four levels, summarised in Table~\ref{tab:strategy}. The
distribution of effort is deliberate: the heaviest unit testing is concentrated
on the two components whose failures would be silent, and the heaviest
system-level effort is concentrated on the deployed endpoints, because a
capability that works locally and not in production has not been delivered.

\begin{xltabular}{\textwidth}{P{2.6cm} P{3.1cm} Y}
\caption{Testing strategy by level.}\label{tab:strategy}\\
\toprule
\bfseries Level & \bfseries Method & \bfseries What it establishes \\
\midrule
\endfirsthead
\toprule
\bfseries Level & \bfseries Method & \bfseries What it establishes \\
\midrule
\endhead
\bottomrule
\endlastfoot
Unit & pytest~\cite{ref:pytest} over synthetic, deterministic inputs & That the
index computation and the feature and supervised builders are correct, including
the leakage properties. \\
Integration & Scheduled pipeline runs against live data and the real feature
store & That ingestion, storage, training, promotion, inference and monitoring
compose correctly under real conditions. \\
System & Direct exercise of every published endpoint on the live deployment & That
the delivered service works from outside, and that unavailable capabilities fail
informatively. \\
Model & Held-out chronological validation, per-horizon metrics, walk-forward
backtesting, baseline comparison & That the forecasts have skill, and how that
skill decays with lead time. \\
\end{xltabular}

\section{Unit Testing}
The unit suite comprises twelve tests. All twelve pass; the run reported
\texttt{12 passed} in 22.01 seconds. The tests are built on synthetic series with
known structure rather than on live data, so they are deterministic and do not
depend on a network.

\subsection{Index Computation Tests}
Seven tests cover the index module.

The first checks the piecewise-linear interpolation at a breakpoint endpoint: a
particulate concentration of 9.0 micrograms per cubic metre sits exactly at the
top of the \emph{good} band and must therefore produce an index of exactly 50.
Testing at an endpoint rather than in the middle of a band is deliberate, since
an off-by-one in the breakpoint tables would be invisible mid-band.

The second checks monotonicity across bands, requiring that a rising sequence of
concentrations produce a non-decreasing sequence of index values, which would
catch a transposed or mis-ordered breakpoint pair.

The third is the most important test in the module, and encodes the design
decision of Section~\ref{sec:impl-aqi}. It supplies a very large carbon monoxide
value alongside clean particulate readings and requires the computed index to be
unaffected - that is, it asserts that gaseous pollutants are excluded from the
computation. Without this test, a later change that added gases back without the
necessary unit conversion would silently corrupt every affected row, because the
index is a maximum over sub-indices and one inflated sub-index dominates.

The remaining four confirm that the category boundaries are assigned correctly at
50, 51 and 301; that the six categories are contiguous, each band beginning
exactly one above the previous band's end, so that no value can fall between
categories; that the hazard predicate fires above the threshold and not below it;
and that a missing concentration yields a missing sub-index rather than an
exception, so that a gap in the provider's data cannot fail a batch.

\subsection{Feature and Supervised-Set Tests}
Five tests cover the feature layer, and two of them are the leakage guards.

The first confirms that the engineered table contains the expected feature
families - a cyclical hour encoding, the target, a change rate, a 24-hour lag, a
24-hour rolling mean, both wind components and the primary key.

The second confirms that the derived primary key is both strictly increasing and
unique within a city, which is the precondition for the upsert semantics on which
every re-run of every pipeline depends.

The third is the cross-city leakage guard. Two synthetic cities are concatenated
and passed through the multi-city builder, and the test requires that
\emph{each} city independently show exactly 24 missing values in its 24-hour lag
column. This can only hold if the two cities were engineered as separate groups;
had the lag been computed over the concatenated frame, the second city would show
none, because it would have silently borrowed the first city's history.

The fourth confirms that the supervised builder emits every declared model
feature and the target, and that the horizon column takes exactly the requested
values - verifying that the horizon really is a feature rather than an implicit
property of the dataset.

The fifth is the temporal leakage guard, asserting that the anchor value attached
to a row at horizon $h$ is the value observed $h$ steps earlier, and that no
anchor is missing after the incomplete rows have been dropped.

\subsection{Coverage Assessment}
The suite covers the two components whose failure modes are silent. It does not
cover the promotion gate, the interval construction, the drift computation or the
API handlers, which were verified by integration and system testing rather than
by unit tests. Section~\ref{sec:testing-threats} treats this as a limitation, and
Chapter~\ref{ch:conclusion} lists it as future work.

\section{Integration and System Testing}
Integration testing was performed by running the pipelines end to end on the
scheduled runners against the real feature store and the real provider. Four
integration properties were confirmed in this way, and each of them had failed at
least once before it held.

\emph{Idempotent ingestion.} The hourly pipeline was run repeatedly within the
same hour, with the deliberate seven-day overlap, and the row count in the feature
store grew only by genuinely new hours, confirming that the composite-key upsert
behaves as designed.

\emph{Resumable backfill.} The backfill was interrupted and restarted, and
resumed only those cities whose stored history did not yet reach the requested
start date, confirming the history-aware resume of
Figure~\ref{fig:code-resume}.

\emph{Registry round trip.} A model trained on one ephemeral runner was loaded by
a different job on a different runner, confirming that the champion genuinely
survives the machine that produced it - the property on which the entire
serverless architecture rests.

\emph{Promotion gating under repetition.} Training was run twice in close
succession on effectively unchanged data. The first run promoted; the second
produced an identical validation RMSE and was refused. This is reported in full in
Section~\ref{sec:testing-gate}.

System testing exercised the deployed public API directly.
Table~\ref{tab:endpoints} reports the result of calling every endpoint against the
live deployment.

\begin{xltabular}{\textwidth}{P{5.4cm} Q{1.3cm} Q{1.6cm} Y}
\caption{System validation against the live deployment.}\label{tab:endpoints}\\
\toprule
\bfseries Endpoint & \bfseries Status & \bfseries Latency & \bfseries Observation \\
\midrule
\endfirsthead
\toprule
\bfseries Endpoint & \bfseries Status & \bfseries Latency & \bfseries Observation \\
\midrule
\endhead
\bottomrule
\endlastfoot
GET /api/health & 200 & 1.32\,s & Reports a forecast document present. \\
GET /api/categories & 200 & 0.82\,s & Six categories with bounds and advice. \\
GET /api/cities & 200 & 0.72\,s & All 22 cities with province and coordinates. \\
GET /api/predictions & 200 & 1.51\,s & 190\,kB payload covering every city. \\
GET /api/predictions/Lahore & 200 & 0.99\,s & 72 hourly points and three daily
points. \\
GET /api/leaderboard & 200 & 0.88\,s & Champion plus three run records. \\
GET /api/evaluation & 200 & 0.74\,s & Ten per-horizon rows and five backtest
folds. \\
GET /api/monitoring & 200 & 0.74\,s & Drift table and forecast-error status. \\
GET /api/explain/Lahore & 200 & 0.75\,s & Sentence plus five ranked
contributions. \\
GET /api/whatif/defaults & 200 & 0.95\,s & Five sliders initialised to live
values. \\
GET /api/predict (arbitrary point) & 200 & 0.98\,s & On-demand forecast for a city
outside the supported list, confirming FR-22. \\
POST /api/whatif & 200 & 1.15\,s & Baseline 38, scenario 39 under a raised wind
speed; both categories returned. \\
POST /api/chat & 503 & 0.74\,s & Returns the documented message that the advisor
credential is not configured on the host. \\
\end{xltabular}

Twelve of the thirteen endpoints serve their documented payload. The thirteenth
behaves exactly as specified when its optional credential is absent: it returns a
specific, actionable message rather than failing inside the client library. This
is the intended behaviour of the exception path in use case UC-7, but it is also
an honest limitation of the current deployment, and it is recorded as such in
Section~\ref{sec:testing-threats}.

\section{Functional Test Cases}
Table~\ref{tab:testcases} records the functional test cases, the requirement each
exercises and its result.

\begin{xltabular}{\textwidth}{Q{1.3cm} P{4.4cm} Q{1.5cm} Y}
\caption{Functional test cases and results.}\label{tab:testcases}\\
\toprule
\bfseries ID & \bfseries Test & \bfseries Req. & \bfseries Result \\
\midrule
\endfirsthead
\toprule
\bfseries ID & \bfseries Test & \bfseries Req. & \bfseries Result \\
\midrule
\endhead
\bottomrule
\endlastfoot
TC-1 & Ingest all cities and confirm 21 raw columns per city with retries
configured. & FR-1 & Pass \\
TC-2 & Engineer features and confirm the 74-column table with cyclical, lag,
rolling and wind-vector families. & FR-2 & Pass \\
TC-3 & Re-run ingestion within the overlap window and confirm no duplicate
rows. & FR-3 & Pass \\
TC-4 & Interrupt and restart the backfill; confirm only incomplete cities are
reprocessed. & FR-3 & Pass \\
TC-5 & Build the supervised set and confirm the horizon column takes exactly the
ten declared values. & FR-4 & Pass \\
TC-6 & Confirm the validation window lies entirely after the training window. &
FR-5 & Pass \\
TC-7 & Confirm all four candidate families are scored in a run. & FR-6 & Pass \\
TC-8 & Retrain on unchanged data and confirm the tie is refused. & FR-7 &
Pass \\
TC-9 & Load the champion on a runner that did not train it. & FR-8 & Pass \\
TC-10 & Confirm 72 hourly points and three daily points for every city. & FR-9 &
Pass \\
TC-11 & Confirm interval bounds present, ordered and horizon-dependent. & FR-10 &
Pass \\
TC-12 & Confirm every forecast point carries a health category. & FR-11 & Pass \\
TC-13 & Confirm a hazard alert with peak, peak hour and first crossing for a city
above the threshold. & FR-12 & Pass \\
TC-14 & Confirm an explanation sentence and five ranked contributions per city. &
FR-13 & Pass \\
TC-15 & Confirm the global importance ranking is produced. & FR-14 & Pass \\
TC-16 & Confirm slider defaults reflect live conditions. & FR-15 & Pass \\
TC-17 & Run a scenario and confirm baseline, scenario, categories and delta. &
FR-16 & Pass \\
TC-18 & Confirm PSI is computed for all seven monitored features with a status
each. & FR-17 & Pass \\
TC-19 & Confirm forecasts are logged and the error report distinguishes ``waiting
for actuals'' from a real score. & FR-18 & Pass \\
TC-20 & Load all four dashboard views against the live backend. & FR-19 & Pass \\
TC-21 & Confirm the four advisor tools are declared and dispatchable. & FR-20 &
Pass \\
TC-22 & Confirm the protocol server exposes the same four tools. & FR-21 &
Pass \\
TC-23 & Request an on-demand forecast for a city outside the supported list. &
FR-22 & Pass \\
TC-24 & Confirm the hourly and daily workflows run to completion on schedule. &
FR-23 & Pass \\
TC-25 & Request the advisor with no credential configured. & NFR-11 & Pass -
returns the documented 503 message. \\
\end{xltabular}

\section{Model Evaluation}

\subsection{Candidate Comparison}
Table~\ref{tab:candidates} reports the champion training run. The supervised set
was 200,000 rows, split chronologically into 160,002 training and 39,998
validation rows.

\begin{xltabular}{\textwidth}{P{5.0cm} C C C Y}
\caption{Candidate comparison on the held-out chronological validation
window.}\label{tab:candidates}\\
\toprule
\bfseries Model & \bfseries RMSE & \bfseries MAE & \bfseries $R^2$ & \bfseries Note \\
\midrule
\endfirsthead
\toprule
\bfseries Model & \bfseries RMSE & \bfseries MAE & \bfseries $R^2$ & \bfseries Note \\
\midrule
\endhead
\bottomrule
\endlastfoot
XGBoost & 19.69 & 12.80 & 0.850 & Promoted champion \\
Random forest & 20.38 & 13.56 & 0.839 & \\
Ridge regression & 22.85 & 16.01 & 0.797 & Linear reference \\
Persistence (baseline) & 25.27 & 15.71 & 0.752 & Not eligible for promotion \\
\end{xltabular}

Three observations follow. The champion reduces RMSE by 22.1\% relative to
persistence, which establishes that the system has real forecasting skill rather
than merely reproducing the current value. The gap between the linear model and
the two ensembles - 22.85 against 20.38 and 19.69 - confirms the reading of the
exploratory analysis in Section~\ref{sec:design-eda}, that the relationship is
genuinely non-linear and interactive rather than an additive combination of
weather variables. And the ordering of the two ensembles is consistent with the
tabular-learning literature~\cite{ref:whytrees}, with boosting ahead of bagging on
a problem of this shape.

One detail is worth drawing out because it demonstrates why RMSE rather than MAE
governs promotion. The persistence baseline has a \emph{lower} MAE than the Ridge
model (15.71 against 16.01) while having a substantially higher RMSE (25.27
against 22.85). Persistence is typically close but occasionally very wrong, which
is exactly the behaviour RMSE penalises and MAE does not; for a health warning
system, the large error during a rapid deterioration is the one that matters, so
the squared metric is the appropriate gate.

\subsection{Accuracy by Forecast Horizon}
Table~\ref{tab:horizon} and Figure~\ref{fig:err-horizon} report accuracy at each
modelled lead time, against the persistence baseline evaluated under the same
mask.

\begin{xltabular}{\textwidth}{C C C C C Y}
\caption{Accuracy by forecast horizon, with the persistence baseline at the same
horizon.}\label{tab:horizon}\\
\toprule
\bfseries Horizon (h) & \bfseries RMSE & \bfseries MAE & \bfseries $R^2$ &
\bfseries Baseline RMSE & \bfseries Improvement \\
\midrule
\endfirsthead
\toprule
\bfseries Horizon (h) & \bfseries RMSE & \bfseries MAE & \bfseries $R^2$ &
\bfseries Baseline RMSE & \bfseries Improvement \\
\midrule
\endhead
\bottomrule
\endlastfoot
1 & 6.10 & 4.01 & 0.985 & 4.54 & $-34.4\%$ \\
2 & 6.92 & 4.44 & 0.981 & 8.43 & $17.9\%$ \\
3 & 7.88 & 5.13 & 0.975 & 11.52 & $31.6\%$ \\
6 & 10.61 & 7.06 & 0.957 & 17.88 & $40.7\%$ \\
12 & 14.78 & 10.46 & 0.917 & 22.89 & $35.4\%$ \\
24 & 21.72 & 15.85 & 0.819 & 25.94 & $16.3\%$ \\
36 & 25.76 & 19.15 & 0.746 & 33.35 & $22.8\%$ \\
48 & 27.56 & 20.64 & 0.708 & 33.07 & $16.7\%$ \\
60 & 27.57 & 20.53 & 0.721 & 35.64 & $22.6\%$ \\
72 & 27.55 & 21.16 & 0.695 & 34.61 & $20.4\%$ \\
\end{xltabular}

\fig{error-by-horizon}{0.86}{RMSE against forecast horizon, compared with the
persistence baseline.}{fig:err-horizon}

The table is the single most informative result in this thesis, and it says three
things.

\emph{Skill decays with lead time, steeply and then gently.} $R^2$ falls from
0.985 at one hour to 0.695 at 72, and RMSE rises from 6.10 to 27.55. Almost all of
the degradation occurs in the first 36 hours; between 36 and 72 hours the error is
essentially flat at 25.8 to 27.6. This plateau is interpretable: beyond about a
day and a half the current pollution state has ceased to carry information, and
the forecast is being driven almost entirely by the forecast weather and the
calendar, whose informativeness does not decay much further over the remaining
window.

\emph{The model beats persistence at every horizon from two hours onward}, by
between 16\% and 41\%, with the largest margins in the 6-to-12-hour range where
the anchor state and the weather forecast are both informative. This is the
result that justifies the system existing.

\emph{At one hour, persistence wins.} Its RMSE of 4.54 is better than the
model's 6.10. This is reported rather than omitted, and it is not a defect: at a
lead time of one hour the air quality is overwhelmingly determined by the air
quality of the preceding hour, and a global model trained jointly across ten
horizons cannot specialise on that regime as tightly as the trivial predictor
does. It is the price of the single-artefact design of
Section~\ref{sec:design-mh}, it is confined to the horizon at which forecasting is
least useful - nobody plans around the next sixty minutes - and the design would
readily accommodate a blend towards persistence at very short lead times if that
were wanted. Reporting it is exactly the kind of disclosure that a single averaged
accuracy figure would have concealed.

\subsection{Walk-Forward Backtest}
Table~\ref{tab:backtest} and Figure~\ref{fig:walkforward} report the five-fold
rolling backtest, in which each fold trains on an expanding window of the past and
tests on the window that follows.

\begin{xltabular}{\textwidth}{C C C C C Y}
\caption{Five-fold walk-forward backtest.}\label{tab:backtest}\\
\toprule
\bfseries Fold & \bfseries Train rows & \bfseries Test rows & \bfseries Test from
& \bfseries RMSE & \bfseries $R^2$ \\
\midrule
\endfirsthead
\toprule
\bfseries Fold & \bfseries Train rows & \bfseries Test rows & \bfseries Test from
& \bfseries RMSE & \bfseries $R^2$ \\
\midrule
\endhead
\bottomrule
\endlastfoot
1 & 33,335 & 33,333 & 2023-08-13 & 21.06 & 0.795 \\
2 & 66,668 & 33,333 & 2024-03-22 & 17.56 & 0.734 \\
3 & 100,001 & 33,333 & 2024-10-30 & 26.44 & 0.778 \\
4 & 133,334 & 33,333 & 2025-06-07 & 19.03 & 0.847 \\
5 & 166,667 & 33,333 & 2026-01-14 & 18.67 & 0.814 \\
\end{xltabular}

\fig{walk-forward}{0.86}{RMSE by walk-forward fold.}{fig:walkforward}

The mean backtest RMSE is 20.55 with a standard deviation of 3.53. Two features of
this result matter more than the mean.

First, the backtest mean is slightly \emph{worse} than the single-split validation
figure of 19.69, and that is the expected and desirable direction. Early folds
train on as little as a fifth of the data, so the backtest deliberately measures
performance under conditions harsher than production; a backtest that flattered the
single split would be evidence of a leak rather than of quality.

Second, the fold-to-fold variation is itself informative. Fold~3, which tests from
30~October~2024, is much the worst at 26.44, while fold~2, testing from late March,
is the best at 17.56. The two are the winter smog season and the clean spring
respectively. The model is therefore measurably harder to fit in the season that
matters most - which is precisely why the nightly retraining and the drift
monitoring exist, and it is a caution against quoting the annual mean error as
though it applied uniformly across the year.

\subsection{The Sequence-Model Comparison}
The LSTM was evaluated on a deliberately fair footing: the same 24-hour-ahead
task, the same chronological split, the same metrics and the same baseline.
Table~\ref{tab:lstm} reports the comparison.

\begin{xltabular}{\textwidth}{P{5.6cm} C C C Y}
\caption{Sequence model against the tree ensemble at a 24-hour
horizon.}\label{tab:lstm}\\
\toprule
\bfseries Model (predict AQI at +24\,h) & \bfseries RMSE & \bfseries MAE &
\bfseries $R^2$ & \bfseries Note \\
\midrule
\endfirsthead
\toprule
\bfseries Model (predict AQI at +24\,h) & \bfseries RMSE & \bfseries MAE &
\bfseries $R^2$ & \bfseries Note \\
\midrule
\endhead
\bottomrule
\endlastfoot
LSTM (two recurrent layers, dense head) & 22.12 & 14.32 & 0.799 & Best at this
single horizon \\
XGBoost (same 24-hour task) & 22.63 & 14.48 & 0.789 & \\
Persistence (baseline) & 28.79 & 17.21 & 0.660 & \\
\end{xltabular}

The sequence model wins, narrowly: 22.12 against 22.63, a 2.3\% reduction in RMSE.
Both learned models beat persistence by roughly 22\%, which independently confirms
that the task has learnable structure.

It was nonetheless not promoted, for three reasons that are worth stating
explicitly because the decision runs against the narrow measurement.

First, \emph{coverage}. The LSTM as constructed predicts one horizon; the
production model serves all ten, at an overall RMSE of 19.69. Replacing the
champion with the sequence model would mean either training ten of them or serving
only one lead time, and the resulting system would be worse on the metric that
actually matters to a user, which is the quality of the whole three-day curve.

Second, \emph{cost}. The sequence model takes roughly 25 minutes to train against
roughly one minute for the ensemble on the same hardware. Within the scheduler's
time limit, ten of them do not fit. The nightly retrain - which is the mechanism
by which the entire system stays current - would have to be abandoned to adopt
the model that is 2.3\% better at one horizon.

Third, \emph{explainability}. Exact Shapley attribution is available in polynomial
time for the tree ensemble~\cite{ref:treeshap} and is not for the recurrent
network, so promoting it would mean withdrawing the explanation that
Chapter~\ref{ch:ui} shows on every forecast.

The finding is reported as a positive one. The sequence model is a legitimate and
competitive alternative, and the comparison is exactly what the tabular-learning
literature~\cite{ref:tabnotall,ref:whytrees} would predict: no decisive advantage
for the deep architecture on a tabular, weather-driven problem. The production
choice was made on the total system objective rather than on a single number, and
the measurement is what allows that to be said rather than assumed.

\subsection{Feature Importance}
Figure~\ref{fig:shap} shows the global feature importance obtained by averaging
absolute Shapley values over a sample of the supervised set.

\fig{shap-importance}{0.78}{Global feature importance by mean absolute Shapley
value.}{fig:shap}

The ranking is physically coherent, which is itself a form of validation. The
anchor state - the current index and its recent means - dominates, as it must
for a persistent quantity. The horizon feature ranks highly, which is direct
evidence that the model has learned to modulate its behaviour by lead time rather
than treating all horizons alike; this is the mechanism by which one artefact can
serve ten horizons. Among the meteorological drivers, the dispersion-related
variables carry the most weight, consistent with the physical account of the smog
season given in Chapter~\ref{ch:literature}. Nothing implausible appears near the
top, which is a check that would have caught a leaked or mis-specified feature.

\section{Uncertainty Calibration}
Prediction intervals are the tenth and ninetieth residual percentiles computed per
horizon on the held-out validation window, giving a nominal 80\% interval. By
construction the intervals attain approximately 80\% coverage on that window, and
they widen with horizon in proportion to the model's actual errors, which is the
behaviour the design intended and which is visible as the widening band in
Figure~\ref{fig:ui-forecast}.

Two limitations must be stated. The construction inherits the exchangeability
assumption of split conformal prediction~\cite{ref:leiconformal}, which a
seasonal, non-stationary series does not strictly satisfy; coverage during a
regime the calibration window did not contain - an unusually severe smog episode,
for instance - should be expected to be worse than nominal. And coverage has not
yet been measured prospectively on live forecasts, because doing so requires a
sufficient volume of matured forecast-outcome pairs. The machinery to perform that
measurement exists in the forecast log described in
Section~\ref{sec:impl-monitoring}, and it is identified as future work.

\section{Operational Evidence from the Running System}

\subsection{The Promotion Gate in Production}
\label{sec:testing-gate}
The leaderboard is the record of the gate's behaviour, and is reproduced in
Table~\ref{tab:leaderboard}.

\begin{xltabular}{\textwidth}{C P{3.0cm} P{2.0cm} C C Y}
\caption{The production model leaderboard.}\label{tab:leaderboard}\\
\toprule
\bfseries Ver & \bfseries Trained & \bfseries Model & \bfseries RMSE &
\bfseries $R^2$ & \bfseries Outcome \\
\midrule
\endfirsthead
\toprule
\bfseries Ver & \bfseries Trained & \bfseries Model & \bfseries RMSE &
\bfseries $R^2$ & \bfseries Outcome \\
\midrule
\endhead
\bottomrule
\endlastfoot
1 & 2026-08-19 17:32 & XGBoost & 20.597 & 0.824 & Promoted - first champion \\
2 & 2026-08-28 14:28 & XGBoost & 19.687 & 0.850 & Promoted - current champion \\
3 & 2026-08-28 15:28 & XGBoost & 19.687 & 0.850 & Rejected - no improvement \\
\end{xltabular}

The three rows demonstrate all three behaviours the gate is required to exhibit.
Version~1 was promoted unconditionally as the first champion. Version~2, trained
after the full historical backfill had completed, improved RMSE from 20.597 to
19.687 - a 4.4\% reduction obtained purely from more history - and displaced the
incumbent, which is the retraining loop delivering measurable value. Version~3,
trained an hour later on effectively unchanged data, produced an identical 19.687
and was refused, because the promotion rule requires strict improvement. That
refusal is the most important row in the table: it is the system declining to
replace its production model for no gain, which is precisely what the gate exists
to do and what an ungated nightly retrain would not have done.

\subsection{Drift Monitoring in Production}
Table~\ref{tab:drift} reports the drift status computed over the most recent
fourteen days against the historical reference.

\begin{xltabular}{\textwidth}{P{4.4cm} C P{2.6cm} Y}
\caption{Population Stability Index by feature, recent fourteen days against the
historical reference.}\label{tab:drift}\\
\toprule
\bfseries Feature & \bfseries PSI & \bfseries Status & \bfseries Interpretation \\
\midrule
\endfirsthead
\toprule
\bfseries Feature & \bfseries PSI & \bfseries Status & \bfseries Interpretation \\
\midrule
\endhead
\bottomrule
\endlastfoot
temperature\_2m & 3.942 & significant & Seasonal transition; the recent window
occupies a temperature range the annual reference visits only briefly. \\
surface\_pressure & 1.253 & significant & Moves with the same seasonal cycle. \\
relative\_humidity\_2m & 0.393 & significant & Monsoon-season humidity against an
all-year reference. \\
aqi & 0.303 & significant & The target distribution has shifted with the season. \\
pm2\_5 & 0.257 & significant & Just above the threshold, consistent with the
index. \\
pm10 & 0.031 & stable & Coarse particulate matter is less seasonally
modulated. \\
wind\_speed\_10m & 0.014 & stable & Wind distribution is broadly stationary. \\
\end{xltabular}

The overall status is \emph{significant}, driven by temperature. This is the
expected result and not a fault, and reading it correctly is the point of the
design in Section~\ref{sec:impl-monitoring}. A fourteen-day window compared
against a multi-year reference will always look shifted on any strongly seasonal
variable, because two weeks of one season cannot resemble an average over all
seasons. That five of the seven monitored features register as significant while
the two least seasonal remain stable is itself corroboration that the measure is
detecting seasonality rather than a data fault.

The correct operational response is the one the system already takes: retrain
nightly, so that the training distribution tracks the current one. What would
constitute a genuine alarm is drift accompanied by rising realised error, which is
why the second monitoring signal exists. At the time of writing that signal
correctly reports that forecasts have been logged and are awaiting matured
observations, rather than displaying a fabricated score.

\section{Non-Functional Testing}
Table~\ref{tab:nfrtest} records how the non-functional requirements were
verified.

\begin{xltabular}{\textwidth}{Q{1.5cm} P{3.4cm} Y}
\caption{Verification of the non-functional requirements.}\label{tab:nfrtest}\\
\toprule
\bfseries ID & \bfseries Requirement & \bfseries Evidence \\
\midrule
\endfirsthead
\toprule
\bfseries ID & \bfseries Requirement & \bfseries Evidence \\
\midrule
\endhead
\bottomrule
\endlastfoot
NFR-1 & Free-tier operation & Ingestion is keyless and free; scheduled compute,
feature storage and static hosting are within free allowances; only the container
host and the optional advisor carry any cost. \\
NFR-2 & Statelessness & A model trained on one runner was loaded by a different
job on another, and the deployed backend serves a champion it never trained. \\
NFR-3 & Performance & All eleven read endpoints responded in 0.72 to 1.51 seconds,
serving pre-computed documents; no page view invokes a model. \\
NFR-4 & Explainability & Every forecast point carries interval bounds; every city
in the live payload carries an explanation object. \\
NFR-5 & Groundedness & The advisor answers only through the four declared tools,
which read the same forecast document the dashboard reads. \\
NFR-6 & Honest framing & The what-if panel states that it is a model simulation
and not causal proof. \\
NFR-7 & Reproducibility & Seeds fixed for every estimator and for subsampling; the
split is defined by time rather than by chance. \\
NFR-8 & Resilience & Explanation, monitoring and forecast logging are each
wrapped; a per-city failure in inference costs one city, not the run. \\
NFR-9 & Auditability & The leaderboard, evaluation report and monitoring report
are committed each night and served publicly. \\
NFR-10 & Portability & The full pipeline runs locally against the Parquet backend
on a platform where the managed client cannot be installed. \\
NFR-11 & Graceful degradation & The advisor returns a specific 503 message when
unconfigured; the what-if tab is hidden in static mode; the monitoring endpoint
has three levels of fallback. \\
NFR-12 & Maintainability & Layered modules with one-directional dependencies;
cities, paths and horizons defined in exactly one module. \\
NFR-13 & Usability & EPA colour scale and category names used throughout; layout
stacks on a narrow viewport; the generation time is displayed in the header. \\
NFR-14 & Politeness to the source & Retry with exponential backoff; three-month
chunking; coverage-based resume avoids refetching. \\
NFR-15 & Security & No credential in the repository; secrets injected from the
environment locally and from the workflow secret store in automation. \\
\end{xltabular}

\section{Threats to Validity}
\label{sec:testing-threats}
Seven threats to the validity of these results are acknowledged.

\textbf{Modelled rather than measured ground truth.} The target is a modelled
index derived from an atmospheric composition product, not a reading from a
reference-grade station in each city. The system therefore demonstrates skill at
forecasting that product. Where the product diverges from a local station, the
forecast inherits the divergence. Validating against ground stations is the single
most valuable extension available to this work.

\textbf{Optimistic weather at training time.} The model is trained on the weather
that actually occurred at the target hour and deployed on the forecast weather for
that hour. The reported validation errors are therefore a mild under-estimate of
operational error, and the gap widens with horizon, since weather forecasts
themselves degrade. This is standard practice, but it is a real bias in the
direction of flattery.

\textbf{Coverage has not been verified prospectively.} The 80\% intervals attain
their nominal coverage on the calibration window by construction; whether they do
so on live forecasts has not yet been measured, and the exchangeability assumption
gives reason to expect degradation during unusual regimes.

\textbf{Limited unit-test coverage.} The suite covers the index computation and
the feature layer thoroughly and leaves the promotion gate, the interval
construction, the drift computation and the API handlers to integration and system
testing. Those components are exercised in production, but a regression in one of
them would not be caught before deployment.

\textbf{A short operational record.} The leaderboard contains three runs and the
realised-error monitor has not yet accumulated matured forecast-outcome pairs. The
gate's behaviour is demonstrated but its long-run behaviour - how often a genuine
improvement occurs, whether a degraded model is ever correctly refused - is not
yet evidenced.

\textbf{A single evaluation regime.} The champion metrics come from one
chronological split whose validation window falls in a particular part of the
year. The walk-forward backtest mitigates this, and its fold-to-fold spread of
17.56 to 26.44 quantifies how much the answer depends on which season is being
tested.

\textbf{No user study.} The interface's claims - that the explanation is
understandable, that the colour scale communicates risk, that the alert prompts
action - are argued from design principles and are not empirically established.

\section{Discussion}
Taken together, the results support four conclusions.

The system forecasts with genuine skill. It reduces RMSE by 22\% against
persistence, and beats it at every horizon from two hours to three days. The skill
is real, it is measured against the correct reference, and it is reported at the
granularity at which it is used.

The skill is honestly bounded. $R^2$ falls from 0.985 to 0.695 across the
forecast window, the backtest is slightly worse than the single split, the winter
fold is markedly worse than the spring fold, and persistence beats the model at a
one-hour lead. None of this is concealed, and the fact that a public dashboard
displays the decaying error curve is a deliberate design position rather than an
oversight.

The operational machinery works and has already earned its place. Retraining after
the backfill improved RMSE by 4.4\%; the gate then refused an identical model an
hour later. Both events are recorded in a leaderboard that a member of the public
can read. This is the difference between a model and a system.

Finally, the model choice was made by measurement rather than fashion. A sequence
model was built, evaluated fairly, found to be 2.3\% better at one horizon, and
declined on grounds of coverage, cost and explainability - with the reasoning
recorded so that a reader can disagree with it.

\section{Summary of Testing and Evaluation}
This chapter reported twelve passing unit tests concentrated on the two
components whose failures would be silent, four integration properties confirmed
on live infrastructure, and thirteen endpoints exercised against the live
deployment, twelve serving their documented payload and the thirteenth degrading
exactly as specified. It reported the champion's accuracy - RMSE 19.69, MAE
12.80, $R^2$ 0.850 - against three alternatives, per-horizon accuracy against a
per-horizon baseline, a five-fold walk-forward backtest averaging 20.55, a fair
comparison with a recurrent sequence model, the global feature importance, the
production leaderboard including a refused promotion, and the current drift
status with its seasonal interpretation. It verified every non-functional
requirement and set out seven threats to the validity of these results. The next
chapter concludes.
